\setstretch{1.15}

\chapter{Related Work}
\label{chapter:related}

This chapter reviews work related to the dissertation's three contributions: storage-backed managed heaps, garbage-collection and data-placement techniques, and dynamic coordination of JVM heap capacity with the operating-system page cache.
The first two contributions extend managed-object capacity onto storage while addressing collector costs under different collection goals.
The third contribution studies how the JVM heap and file cache should share DRAM, and can complement storage-backed heaps without itself implementing one.

\section{Extending Managed Heaps Beyond DRAM}

Several systems increase the memory available to managed applications by using secondary storage, non-volatile memory (NVM), or remote memory.
Their designs differ in whether objects retain their managed representation, how the collector discovers and reclaims them, and whether the target is local block-addressable storage or another memory tier.

\subsection{Local Storage-Backed Heaps}

LOOM extends object memory across DRAM and secondary storage~\cite{kaehler:loom}.
Other systems, including Abdelrahman and Pande's design, Melt, and LeakSurvivor, move cold or stale reachable objects to disk and restore them when accessed~\cite{abdelrahman:memorygc,bond:melt,leaksurvivor}.
These designs target objects with infrequent accesses and use mechanisms that differ from a memory-mapped managed heap.
Melt and LeakSurvivor maintain object-location metadata in the JVM, while LeakSurvivor also defers reclamation on disk until space is needed.
TMO monitors DRAM pressure and moves cold pages to slower tiers, including NVMe SSDs, below the JVM's object model~\cite{weiner:tmo}.

TeraHeap and Chipper instead expose a storage-backed heap in which objects retain their Java representations and are accessed through memory mapping.
TeraHeap integrates this heap with the throughput-oriented Parallel Scavenge collector and avoids scanning storage-resident objects during ordinary collection~\cite{teraheap,kolokasis:teracache}.
Chipper addresses a different constraint: it integrates storage migration with G1's incremental, pause-bounded collection policy, accounting for migration work when selecting regions for evacuation.
These designs share the goal of reducing serialization and storage-resident GC work, but their collector mechanisms and latency objectives differ.

Table~\ref{tab:related} summarizes selected systems that extend managed heaps beyond local DRAM according to developer effort, whether GC scans the secondary tier, and device type.
The table reflects the taxonomy used to compare TeraHeap with prior heap-scaling systems; it is not a complete comparison of low-pause collector designs.

\setlength{\unitlength}{1em}

\begin{table}[t]
\centering
\caption{Comparison of selected systems that extend managed heaps beyond DRAM.}
\resizebox{0.8\columnwidth}{!}{%
\begin{tabular}{lccccc}
\hline
\multicolumn{1}{c}{\multirow{2}{*}{\textbf{System}}} &
  \multirow{2}{*}{\textbf{\begin{tabular}[c]{@{}c@{}}Developer \\ Effort\end{tabular}}} &
  \multicolumn{2}{c}{\textbf{\begin{tabular}[c]{@{}c@{}}GC Scans\\ on the Device\end{tabular}}} &
  \multicolumn{2}{c}{\textbf{Device Type}} \\ \cline{3-6}
\multicolumn{1}{c}{}         &      & \textbf{Lazy} & \textbf{Eager}       & \textbf{Block-addressable} & \textbf{Byte-addressable} \\ \hline
  LeakSurvivor~\cite{leaksurvivor}                  & High & (\checkmark)  & -                    & \checkmark                 & -                         \\
  Melt~\cite{bond:melt}                             & None & (\checkmark)  & -                    & \checkmark                 & -                         \\
  Write Rationing~\cite{akram:ration}               & None & -             & \checkmark           & -                          & \checkmark                \\
  Panthera~\cite{wang:Panthera}                     & Low  & -             & \checkmark           & -                          & \checkmark                \\
  NVM-friendly-GC~\cite{nvm-gc-eurosys}             & None & -             & \checkmark           & -                          & \checkmark                \\
  JPD Heap~\cite{jpd_heap}                     & High & -             & \checkmark           & -                          & \checkmark                \\
  GCMove~\cite{gcmove}                      & None & -             & \checkmark           & -                          & \checkmark                \\
  ThinGC~\cite{thingc}                      & None & -             & \checkmark           & -                          & \checkmark                \\
  TMO~\cite{weiner:tmo}                          & None & -             & \checkmark           & \checkmark                 & -                         \\
  Semeru~\cite{wang:semeru}                       & None & -             & \checkmark           & -                          & \checkmark                \\
  MemLiner~\cite{wang:memliner}                     & None & -             & \checkmark           & -                          & \checkmark                \\
  Mako~\cite{ma:mako}                         & None & -             & \checkmark           & -                          & \checkmark                \\ \hline
\textbf{TeraHeap (this thesis)} & Low  & \checkmark    & -                    & \checkmark                 & \checkmark                \\ \hline
\end{tabular}%
}
\label{tab:related}
\end{table}


\subsection{Byte-Addressable NVM and Remote Memory}

Work on byte-addressable NVM includes systems that place objects according to access or write hotness to reduce writes or improve performance~\cite{akram:ration,akram:crystal,gcmove,thingc}.
NVM-friendly-GC optimizes G1 for an NVM-resident heap and uses DRAM as a cache for collector operations~\cite{Xu:gc_evaluation}.
Panthera and JPDHeap place managed objects across DRAM and NVM, but allow the collector to scan NVM-resident objects~\cite{wang:Panthera,jpd_heap}.
These systems target byte-addressable media, whose access granularity and capacity differ from local NVMe SSDs.

Remote-memory systems address the cost of tracing and collecting objects across a network.
MemLiner improves far-memory locality by aligning collector access order with mutator access patterns~\cite{wang:memliner}.
Semeru separates execution and memory responsibilities across JVM instances, while Mako offloads object scanning and evacuation to memory servers~\cite{wang:semeru,ma:mako}.
DGC offloads concurrent marking to a remote garbage-collection service over RDMA~\cite{hongtao:dgc}.
SmartSweep explores reclamation of remote or tiered managed heaps by returning garbage-rich regions to local memory for collection~\cite{kolokasis:smartsweep}.
SmartSweep is an exploratory direction in this dissertation rather than a fourth principal contribution.

\section{Reducing Managed-Memory and Data-Movement Costs}

\subsection{Reducing Garbage-Collection Cost}

Some systems reduce GC cost while keeping application data in DRAM.
Facade and Gerenuk transform application code to use off-heap representations, reducing the amount of object data traced by the JVM~\cite{nguyen:facade,navasca:gerenuk}.
Broom groups objects with similar lifetimes into regions, while DSA uses programmer annotations to allocate selected data structures separately and improve tracing locality and parallelism~\cite{gog:broom,cohen:dsa}.
Yak allocates objects by programmer-defined epochs, and Deca uses lifetime information to manage Spark objects~\cite{khanh:yak,lu:deca}.
BridgeGC uses framework annotations to avoid repeatedly tracing and copying long-lived data until the framework signals its release~\cite{wang:bridgegc}.
NG2C and allocation-site profiling identify long-lived objects for pretenuring~\cite{bruno:n2gc,pretenuring,Blackburn:2001:PJ}.
These approaches reduce tracing, copying, or allocation overhead, but do not provide the capacity of a storage-backed heap.

\subsection{Reducing Serialization and Data-Movement Cost}

Serialization libraries such as Thrift, Kryo, and Protocol Buffers provide general-purpose formats and implementations for converting object graphs to and from byte streams~\cite{thrift,kryo,protocol-buffers}.
Other systems reduce conversion or transfer costs by moving managed object graphs directly between JVMs, as in Skyway, or by bypassing conventional serialization for RDMA or inter-JVM transfers, as in Naos and ZCOT~\cite{nguyen:skyway,taranov:naos,wu:zero_change}.
SSDStreamer uses a lightweight serializer and a stream-oriented I/O path for SSD-backed caching~\cite{ssdstreamer}.
Specialized hardware and techniques that reduce copies across buffers further accelerate serialization and data transfer~\cite{jang:arch,pourhabibi:optimus,matei:champions,ibanez:zerializer}.
These systems address S/D or copy costs, while TeraHeap and Chipper retain managed object representations in the storage-backed heap and address the collector costs that arise from expanding the heap onto storage.

\section{Heap Sizing and DRAM Coordination}

\subsection{JVM Heap-Resizing Policies}

Prior heap-sizing work uses heuristics, analytical models, profiling, or feedback control to adapt heap capacity to application behavior.
Heuristic policies have been proposed for Boehm-Demers-Weiser GC and JVMs, and other work studies heap resizing under paging, allocation stalls, or changing application requirements~\cite{brecht:nogc,boehm:gc,zhang:pamm,grz:isla_vista,bruno:resize,cook:auto,simao:adaptability,zgc:heap_resize}.
Model-based policies estimate heap needs or paging costs from application and operating-system measurements~\cite{yang:real_mem,sun:multi_space_heap,tay:page_fault,vengerov:gc_analysis,singer:economics,kirisame:membalancer}.
Control-theoretic approaches use feedback to maintain a target GC objective or coordinate memory across database structures~\cite{white:pid,storm:db2}.
Cross-runtime systems allocate memory among multiple JVMs or managed processes to reduce paging or improve aggregate throughput~\cite{alonso:advisor,hertz:waste_not,cameron:forseti,wang:elasticMem}.

These policies address heap capacity, GC overhead, paging, or shared-memory allocation, but generally do not distinguish DRAM used by the JVM heap from DRAM used by the operating system to cache storage-backed data.
That distinction matters when heap and I/O demands change together.

\subsection{Coordinating the Heap with the I/O Cache}

Application-specific cache systems allow developers to align caching with application access patterns~\cite{lee:p2cache}.
Symbiosis coordinates user-space block caches with the operating-system I/O cache, using offline profiling to guide its decisions~\cite{dai:symbiosis}.
FlexHeap instead resizes the JVM heap using feedback about both GC cost and I/O stalls, dynamically reallocating DRAM between the heap and the I/O cache.
It can therefore complement TeraHeap and Chipper, where the OS page cache affects access to storage-backed heap objects, but it is not itself a storage-backed collector.

\section{Positioning of This Dissertation}

The related work exposes three distinct system problems.
TeraHeap increases managed-heap capacity for throughput-oriented applications while limiting collection work over storage.
Chipper extends storage-backed heaps to G1 and accounts for migration cost within incremental, pause-bounded collection.
FlexHeap addresses a complementary resource-allocation problem by balancing JVM heap capacity against the OS page cache according to GC and I/O costs.
Together, these contributions study storage capacity, collector policy, and DRAM allocation as related but separate design dimensions.
